\section{Related Work}
%To solve the problems of connection setup and transport layer security (TLS) handshaking delay for every HTTP request-response flow, Google introduced a new protocol in 2012, called SPDY \cite{Erman,SPDYingupweb} that can multiplex more than one HTTP requests/responses into one TCP connection. Although SPDY reduces the connection establishment time, however, due to TCP’s in-order packet delivery, all the HTTP streams in a multiplexed TCP connection get blocked with SPDY even for a single packet drop event from any one of the HTTP streams, which is known as Head of the Line (HoL) blocking problem. To address HoL blocking, Google has introduced another new experimental protocol called QUIC \cite{Cui2017,Megyesi2016,roskind2015quic} that uses UDP based multiplexed streaming to provide low latency connections for web access. As shown in Fig.~\ref{fig:quic-protocolstack}, QUIC uses UDP as the underlying transport layer protocol to solve the HoL blocking issue with SPDY, and therefore it has its own congestion control and flow control mechanisms \cite{Cui2017,roskind2015quic}. The key benefits of QUIC over TCP+TLS are as follows - (a) reduction in connection establishment latency, (b) multiplexing of multiple application flows without head-of-line blocking, (c) improved congestion control over multiplexed application flows, (d) forward error correction for low latency communication, and (e) connection migration to support seamless mobility in the network.
%
%\begin{figure}[!t]
%	\centering
%	\includegraphics[width=0.7\linewidth]{img/quic-protocolstack}
%	\caption{QUIC Protocol Stack}
%	\label{fig:quic-protocolstack}
%\end{figure}
%In \cite{Cui2017}, Cui \textit{et al.} have shown that QUIC has the potential to change the web and make it faster and more secure. Megyesi \etal \cite{Megyesi2016} have compared QUIC with SPDY and HTTP and found that QUIC outperforms those protocols for web access under low network bandwidth availability. However, they have also noted that QUIC does not work well with the large objects, such as large file transfer. In \cite{Biswal2017}, Biswal \etal have found that QUIC can load web pages much faster than HTTP2.0. They have also shown that QUIC’s performance increases with increasing size of object and number of objects on a page, which is little contrary to the observation given by \cite{Megyesi2016}. %Lychev \etal  \cite{Lychev2015} have analyzed security aspects of QUIC and have presented the first security model for QUIC. They have found that simple bit-flipping and replay attack can prevent QUIC from achieving its goal of low latency. They have also claimed that this weakness is there because of the goal of low-latency. In \cite{bakri2015http}, the authors have developed a service for multi-user virtual world by exploring the low latency features of HTTP2.0 and QUIC.
Various recent studies \cite{Biswal2017,Megyesi2016} have revealed that QUIC can improve the web performance by reducing the page load time even at poor network conditions. Among them, a few works have explored the adaptive streaming performance over QUIC. In~\cite{bhat2017not}, the authors have empirically shown that DASH suffers over QUIC. Although they have given the first indication that the current ABR techniques might not perform well over QUIC, their analysis is mostly focused on buffer-based ABR and does not look into various QoE metrics as explored in the recent literature~\cite{yin2015control,mao2017neural}. In a follow-up work~\cite{bhat2018improving}, the authors have explored the QUIC retransmissions to improve the buffer-based ABR over DASH. In~\cite{van2018empirical}, the authors have used an emulated setup to analyze the buffer-based ABR techniques over QUIC and also proposed a QoE prediction mechanism for adaptive streaming over QUIC. Also, these existing studies have indicated that QUIC might not suit well for buffer-based ABR. They have not explored the performance of advanced ABR techniques over QUIC. We argue that there is a requirement to analyze the recent ABR techniques like MPC~\cite{yin2015control} and Pensieve~\cite{mao2017neural} over QUIC because these recent studies have indicated that buffer-based techniques are aggressive towards video-bitrate maximization whereas suffers in terms of playback smoothness and rebuffering. To the best of our knowledge, this is the first study that explores the performance of advanced ABR techniques over QUIC as the end-to-end transport protocol. 

%In summary, most of these works claim that QUIC has the proficiency to improve the performance of web access. Therefore, it would be interesting to see how web based video streaming services perform when the underlying TCP+TLS stack is replaced by QUIC. However, to the best of our knowledge, no recent work has looked into web based adaptive streaming performance over QUIC.